% Chapter 2 -- Background and Related Work (\input by main.tex)
% Positions this thesis's classical-clustering + binary-classifier + PCN-completion
% pipeline against the deep-learning LiDAR-segmentation and point-completion
% literature, and carries the B3 mandatory literature-positioning table (tab:positioning;
% now Table 2.2, since the completion-metric tab:completion-metrics in 2.1.4 precedes it)
% that defends examiner Q2 ("why F1 0.808 / R 0.73 vs SOTA?") and Q10 ("what is yours
% vs. reimplementation?").
%
% CITATIONS: unlike the results/methods siblings (which cite the internal Finding record),
% this is a literature chapter and uses classic BibTeX. \bibliography{../../references}
% resolves to docs/references.bib (same relative-path pattern as
% docs/report/progress_report_2026_05_12.tex, confirmed to compile in this repo).
%
% Scope decision (settled with advisor/author): RELATED-WORK-FOCUSED. Method mechanics
% (PointNet architecture, HDBSCAN internals, PCN encoder/decoder, Chamfer distance) are
% explained in Chapter~\ref{ch:pipeline} (pipeline) and Section~\ref{ch:completion-method} (completion) and are NOT re-derived
% here; this chapter carries only the light conceptual grounding needed to read the
% related work and the positioning table.
%
% Honesty-phrasing rules applied (personal_agenda P1):
%   - the positioning table (tab:positioning, Table 2.2) is REQUIREMENT-ONLY (approach / segmentation
%     supervision / output type); it deliberately carries NO accuracy column, because
%     mixing mIoU, per-class car IoU, PQ (and its variants), and this work's point-level
%     car F1 in one column invites a comparison those metrics do not support
%     (advisor feedback, 2026-09-07). The accuracy gap is stated in the prose below the
%     table instead.
%   - every quantitative claim about a published method traces to a cited primary source;
%     the mIoU / PQ ranges quoted in the post-table prose are SemanticKITTI test-set
%     figures verified against the paper
%   - this work's headline numbers (F1 0.808 / R 0.730 / P 0.905, seq 08;
%     Chapter~\ref{ch:evaluation}, Finding~#34) appear in the post-table prose with the
%     point-IoU>=0.3 protocol named, NOT in the table
%   - "recall" is stated against its car-only, >=10-surviving-points denominator where used
%
\chapter{Literature Review \& Theoretical Background}
\label{ch:background}

This chapter places the thesis in its research context. It reviews the two bodies of
work the pipeline draws on --- LiDAR point-cloud segmentation and object-level
point-cloud completion --- and then positions the specific design choices of this
work against them. The organising contrast runs throughout: the dominant approach to
LiDAR segmentation is a single end-to-end network trained on dense per-point
annotation, whereas this thesis builds an interpretable pipeline from training-free
geometric clustering plus a lightweight learned classifier, and treats completion as a
separate, downstream problem. Section~\ref{sec:bg-positioning} makes that contrast
quantitative in a positioning table. The chapter is in two parts. Section~\ref{sec:bg-litreview} is the literature review ---
what exists in LiDAR segmentation and object completion, and where this work sits within it
--- and Section~\ref{sec:bg-theory} is a theoretical-background section that fixes the
notation and standard definitions for the operators the pipeline is built from. The
pipeline-specific architectures, configurations, and design choices of each component are
deferred to the method chapter (Section~\ref{ch:pipeline} for detection,
Section~\ref{ch:completion-method} for completion).

\textbf{Scope of the review.} The field around automotive-LiDAR perception is far larger
than any one chapter can survey, so the review is bounded by relevance to the thesis's own
design and evaluation. Included are the deep and classical LiDAR-segmentation paradigms the
detection front-end is built from and positioned against (Sections~\ref{sec:bg-deep}
and~\ref{sec:bg-classical}), and object-level point-cloud completion together with its
evaluation metrics --- the setting of the thesis's contribution (Section~\ref{sec:bg-completion}).
Left deliberately outside are camera--LiDAR fusion and SLAM or map optimisation, both
orthogonal to a single-scan, LiDAR-only object pipeline, and full 3D object-detection
benchmarks scored by bounding-box average precision, which measure a different output than the
point-level clusters this pipeline emits --- an exclusion argued in full where it naturally
arises (Section~\ref{sec:bg-positioning}). The conceptual grounding and notation needed to
read this related work, but no more, are gathered in Section~\ref{sec:bg-theory}.

\section{Literature Review}
\label{sec:bg-litreview}

\subsection{LiDAR Point Clouds and the SemanticKITTI Benchmark}
\label{sec:bg-dataset}

A rotating LiDAR sensor returns, for each scan, an unordered set of 3D points sampled
from the surfaces around the vehicle. The points are sparse, irregularly spaced, and
strongly range-dependent: nearby objects are dense, distant ones may be a handful of
returns, and every object is seen from a single viewpoint and is therefore partial.
These properties --- irregularity, sparsity, single-view partiality --- shape every
method discussed below and every difficulty this thesis reports.

The KITTI Vision Benchmark~\cite{geiger2012we} established the standard automotive
LiDAR setting (a Velodyne HDL-64E on a car, with synchronised camera and pose). Its
segmentation-oriented successor, SemanticKITTI~\cite{Behley_2019_ICCV}, provides
dense per-point semantic labels for the 22 sequences of the KITTI odometry split
(sequences 00--10 are released with labels; the labels for sequences 11--21 are withheld
and scored only through the online test server), and a later extension adds temporally
consistent instance identities for a LiDAR-based panoptic benchmark~\cite{behley2021panoptic}. Three tasks are defined on
it. \emph{Semantic segmentation} assigns every point a class and is scored by the mean
intersection-over-union across classes (mIoU). \emph{Instance segmentation} separates
individual objects within the ``thing'' classes (car, pedestrian, cyclist, and so on).
\emph{Panoptic segmentation} unifies the two and is scored by panoptic quality (PQ),
which factorises into a segmentation-quality term (SQ, the mean IoU of matched
segments) and a recognition-quality term (RQ, an F$_1$-like matching score);
PQ$^{\text{Th}}$ restricts PQ to the ``thing'' classes, and PQ$^\dagger$ is a variant
that scores ``stuff'' by IoU. These definitions matter for
Section~\ref{sec:bg-positioning}, because they are the metrics the literature reports
and none of them is the point-level per-class F$_1$ this thesis uses.

The decisive property of SemanticKITTI for this work is its \emph{annotation cost}.
The labels are dense and per-point across billions of points, and the panoptic
extension adds a manual instance-association pass over sequences. This is the resource
the mainstream methods below consume and the one this thesis is built to avoid: its
segmentation stage requires no labelled point clouds at all, and its classifier is
trained on clusters mined and weakly labelled automatically from the pipeline's own
output (Chapter~\ref{ch:pipeline}).

\subsection{Deep Segmentation of LiDAR Point Clouds}
\label{sec:bg-deep}

The current benchmark leaders are deep networks trained end-to-end on the dense labels.
They differ mainly in how they impose structure on the irregular point set.
\emph{Projection-based} methods rasterise the scan into a 2D range image and apply
image convolutions: SqueezeSeg~\cite{10.1109/ICRA.2018.8462926} and
RangeNet++~\cite{milioto2019rangenet++} are representative, and are fast because they
reuse mature 2D machinery, at the cost of discretisation artefacts at depth
discontinuities. \emph{Point-based} methods operate on the raw points through
permutation-invariant operators: PointNet~\cite{Qi_2017_CVPR} introduced the
shared-MLP-plus-symmetric-pooling design, PointNet++~\cite{NIPS2017_d8bf84be} added
hierarchical local grouping, and RandLA-Net~\cite{Hu_2020_CVPR} made the point-based
route scale to full driving scans through efficient random sampling. \emph{Voxel-based}
methods discretise 3D space and run sparse convolutions; Cylinder3D~\cite{zhu2021cylindrical}
uses a cylindrical partition matched to the sensor's polar geometry and reports
\num{67.8} mIoU on the SemanticKITTI test set, among the higher single-scan results.

Panoptic methods extend semantic backbones with instance reasoning.
Panoptic-PolarNet~\cite{zhou2021panopticpolarnet} predicts instances proposal-free in a
polar bird's-eye-view; DS-Net~\cite{hong2021dsnet} clusters in a learned embedding via a
dynamic shifting module; Panoster~\cite{gasperini2021panoster} learns instance ids
end-to-end; and EfficientLPS~\cite{sirohi2021efficientlps} targets an efficient
architecture for the joint task. The dataset authors' own baselines pair a semantic
network with a 3D detector to supply instances~\cite{behley2021panoptic}.

Two properties are common to this whole family and frame the thesis's alternative.
First, all of them require the dense per-point (and, for panoptic, per-instance)
supervision described above; their accuracy is inseparable from that annotation budget.
Second, they are end-to-end networks whose decisions are not individually inspectable
--- a misclassified region has no stage-level explanation. PointNet and PointNet++ are
relevant to this thesis for a narrower reason as well: they are the backbone family from
which its car/not-car classifier is built (Chapter~\ref{ch:pipeline}). The classifier, however, is a
small binary network applied to already-segmented clusters, not a full-scene
segmentation model, and it is trained without any hand-drawn labels.

\subsection{Classical Geometric Segmentation}
\label{sec:bg-classical}

Before end-to-end learning dominated the benchmarks, LiDAR objects were extracted by
geometric reasoning, and that route is the one this thesis takes for the segmentation
stage. The standard recipe removes the ground plane --- commonly by RANSAC plane
fitting --- and clusters the remaining points into objects. The clustering step is
where the design choice lives. Fixed-radius methods such as Euclidean clustering and
DBSCAN~\cite{ester1996dbscan} group points within a single distance threshold
$\varepsilon$; this is fast and fully interpretable, but a single $\varepsilon$ cannot
serve both the dense near field and the sparse far field of a LiDAR scan.
HDBSCAN~\cite{campello2013hdbscan}, in the widely used implementation
of~\cite{mcinnes2017hdbscan}, replaces the fixed radius with a hierarchy over density
levels and extracts clusters that are stable across scales, which is better matched to
range-varying point density.

This clustering-plus-classification recipe remains an active design point wherever
interpretability and real-time efficiency are valued over benchmark accuracy. Recent
work for autonomous racing, for instance, pairs bird's-eye-view clustering with a
lightweight machine-learning vehicle classifier to reach long-range detection without a
deep dense-feature network~\cite{lim2025longrange}. The pipeline in this thesis is of
the same family --- ground removal, clustering, then a learned per-object classifier ---
applied to the SemanticKITTI setting.

The appeal of this route is that it needs no training data for segmentation and every
stage is inspectable. Its weaknesses are equally concrete and are documented directly in
this thesis: density-based clustering has no notion of objectness, so it fragments a
single car into several clusters when internal density gaps appear, and this
fragmentation --- not classification error --- is the dominant recall limiter
(Section~\ref{sec:recall-analysis}). The positioning below should be read with that trade in view: the
interpretability and near-zero annotation of the classical route are bought at a
measurable recall cost relative to the supervised networks.

\subsection{Object-Level Point-Cloud Completion}
\label{sec:bg-completion}

The second phase of the pipeline reconstructs a denser shape from the sparse, single-view
partial cloud of a detected car. Point-cloud completion learns this partial-to-complete
mapping. PCN~\cite{yuan2019pcnpointcompletionnetwork} established the coarse-to-fine
design --- a PointNet encoder producing a global shape code, decoded first to a coarse
set and then folded to a dense output --- and remains a strong, compact baseline. Later
work improves detail through attention and progressive refinement:
PoinTr~\cite{yu2021pointr} casts completion as set-to-set translation with a
geometry-aware transformer, SnowflakeNet~\cite{xiang2021snowflakenet} generates points
by hierarchical ``snowflake'' deconvolution, and SeedFormer~\cite{zhou2022seedformer}
introduces patch-seed representations. These models are trained on synthetic CAD shapes,
most commonly ShapeNet~\cite{chang2015shapenetinformationrich3dmodel}, with partial
inputs produced by virtual rendering.

Two points from this literature bear directly on later chapters. First, models trained
on clean synthetic partials face a synthetic-to-real domain gap when applied to noisy,
sparse LiDAR clusters; this thesis adopts PCN and confronts that gap head-on
(Section~\ref{ch:completion-method}), and its comparison of PCN against a transformer alternative found no
decisive real-data advantage for the heavier model (Section~\ref{sec:donor-validation}). Second, and more
consequentially, the standard completion metric --- Chamfer distance to a reference
shape --- presupposes a complete ground-truth surface, which real accumulated LiDAR does
not provide: the accumulation is itself one-sided. Section~\ref{sec:donor-validation} shows this makes pseudo
ground-truth Chamfer scoring invalid on real cars and develops an alternative metric in
its place. This should be distinguished from \emph{semantic scene completion}, which
fills a whole voxelised scene from a single scan; the present work completes individual
segmented objects, a different and narrower task.

Because real automotive LiDAR offers no complete target, a small toolkit of
completion metrics has grown up around the KITTI cars. The established quantitative
measures are the \emph{unidirectional} Chamfer and Hausdorff distances (UCD and UHD),
which score a completion by how closely it reproduces the points of its own partial
input: Chen et al.~\cite{chen2020unpaired} apply UHD directly to KITTI cars, and
P2C~\cite{cui2023p2c} uses UCD together with a region-aware variant (RCD) that localises
the comparison around observed skeleton points but still anchors it to seen geometry.
PCN~\cite{yuan2019pcnpointcompletionnetwork} instead scores its KITTI completions by
\emph{Fidelity} (how well the input is preserved in the output), \emph{Minimal Matching
Distance} (Chamfer distance to the nearest complete car in a synthetic reference set), and
\emph{Consistency} (agreement between the completions of one instance across consecutive
frames). Each of these references something other than the real occluded surface --- the
observed input (UCD, UHD, RCD, Fidelity), a synthetic prior (MMD), or the model's own
outputs (Consistency) --- so none directly measures how much genuinely unobserved surface a
completion recovers. Indeed, because a completion that merely copies its input is scored
best by an input-referenced measure, these metrics are used in the unpaired-completion
literature alongside a separate plausibility or adversarial term rather than as a standalone
ranking of recovered surface. Closest to the present work, Ren
et al.~\cite{ren2022selfsupervised} complete vehicles in real traffic scenes and build a
pseudo ground truth by accumulating each tracked instance across consecutive frames, then
score the completed cloud against that accumulation with Chamfer-based distances. That
accumulation is scored \emph{as a whole}, however --- indeed one of its three distance
terms measures input fidelity directly --- so it mixes the near surface the input already
observed with the little new surface the other views add, a completion is not penalised
for merely reproducing its input, and no out-of-box hallucination guard flags surface
invented outside the car. Section~\ref{sec:donor-validation} adopts the
same multi-frame idea but restricts scoring to the surface the input frame did not
see, and adds an explicit hallucination guard, precisely so that reproducing the input
earns nothing.

Table~\ref{tab:completion-metrics} summarises this landscape against the two properties
the completion claim needs: whether a metric credits genuinely never-observed surface ---
and how directly --- and whether it guards against hallucination. The existing measures
range from penalising added surface, through crediting it only indirectly or mixed with
already-seen surface, to not measuring it at all, and none guards hallucination.

\begin{table}[H]
  \centering
  \caption[Completion metrics and recovered occluded surface]{Established real-LiDAR
    completion metrics, by what each scores a completion against, whether it bounds surface
    invented outside the car (``guard''), and how it credits recovered never-observed
    surface (``unseen''): \emph{no} = does not credit it (the input-fidelity measures in
    fact penalise it); \emph{indirect} = credited only through a proxy such as resemblance
    to a synthetic car; \emph{partial} = credited, but mixed with already-seen surface;
    \emph{yes} = scored in isolation. Each measure is valid for its own purpose; none,
    however, isolates recovered unseen surface or guards hallucination, which is what the
    donor metric of Section~\ref{sec:donor-validation} adds.}
  \label{tab:completion-metrics}
  \small
  \renewcommand{\arraystretch}{1.2}
  \setlength{\tabcolsep}{6pt}
  \begin{tabular}{@{}p{2.6cm}p{4.0cm}cc@{}}
    \toprule
    Metric & Scored against & Guard & Unseen \\
    \midrule
    UCD~\cite{cui2023p2c}       & observed input & no & no \\
    UHD~\cite{chen2020unpaired} & observed input & no & no \\
    Fidelity~\cite{yuan2019pcnpointcompletionnetwork} & observed input & no & no \\
    MMD~\cite{yuan2019pcnpointcompletionnetwork} & synthetic CAD prior & no & indirect \\
    Consistency~\cite{yuan2019pcnpointcompletionnetwork} & model's own frames (cross-frame) & no & no \\
    meanCD~\cite{ren2022selfsupervised} & full multi-frame accumulation & no & partial \\
    \midrule
    \textbf{Donor cov.} (ours) & donor-only never-observed surface & yes & yes \\
    \bottomrule
  \end{tabular}
\end{table}

\subsection{Positioning of This Work}
\label{sec:bg-positioning}

The pipeline studied here is a car-only instance detector built from training-free
clustering and a binary classifier, followed by object completion. It occupies a
different point in the design space from the benchmark leaders of
Section~\ref{sec:bg-deep}. Table~\ref{tab:positioning} sets that position out along the
axes on which it actually differs --- approach, segmentation supervision, and output type
--- rather than on an accuracy score. A shared accuracy column is deliberately omitted:
the published methods are scored by mIoU and panoptic quality across all classes, this
work by a car-only point-level F$_1$, and placing those figures in one column would invite
a comparison the metrics do not support. The accuracy gap is real, and is stated plainly
in the text below the table instead.

\begin{table}[H]
  \centering
  \caption[Positioning against representative SemanticKITTI methods]{Positioning
    against representative SemanticKITTI methods by approach, segmentation supervision,
    and output type --- the axes on which this work actually differs. A reported-accuracy
    column is deliberately omitted: the published methods are scored by mIoU and panoptic
    quality across all classes and this work by a car-only point-level F$_1$, so no single
    accuracy column would be comparable across rows. The accuracy gap is stated in the
    text.}
  \label{tab:positioning}
  \small
  \renewcommand{\arraystretch}{1.25}
  \setlength{\tabcolsep}{5pt}
  \begin{tabular}{@{}p{2.7cm}p{3.7cm}p{3.5cm}p{3.0cm}@{}}
    \toprule
    Method & Approach & Supervision & Output \\
    \midrule
    RangeNet++~\cite{milioto2019rangenet++} & Range-image CNN, semantic & Dense per-point & Per-point class \\
    RandLA-Net~\cite{Hu_2020_CVPR} & Point-based, semantic & Dense per-point & Per-point class \\
    Cylinder3D~\cite{zhu2021cylindrical} & Voxel (cylindrical), semantic & Dense per-point & Per-point class \\
    \addlinespace[2pt]
    RangeNet++ \& PointPillars~\cite{behley2021panoptic} & Semantic net + detector, panoptic & Dense per-point + boxes & Per-point class + instance \\
    Panoptic-PolarNet~\cite{zhou2021panopticpolarnet} & End-to-end panoptic & Dense per-point + instance & Per-point class + instance \\
    DS-Net~\cite{hong2021dsnet} & End-to-end panoptic & Dense per-point + instance & Per-point class + instance \\
    EfficientLPS~\cite{sirohi2021efficientlps} & End-to-end panoptic & Dense per-point + instance & Per-point class + instance \\
    \midrule
    \textbf{This work} & Geometric clustering + binary classifier (+ completion) & None for segmentation; auto-mined weak labels for the classifier & Point-level car clusters (+ completion) \\
    \bottomrule
  \end{tabular}
\end{table}

Read as requirements, the table states the thesis's position plainly. On their own
benchmarks the supervised networks reach \numrange{52.2}{67.8} mIoU (semantic) and
\numrange{37.1}{57.4} PQ (panoptic) over the official test sequences~11--21, scored
across all classes; this work reports a car-only point-level F$_1$ of \num{0.808}
(recall \num{0.730} against the $\ge$-10-surviving-points denominator of
Chapter~\ref{ch:evaluation}, precision \num{0.905}) on sequence~08. These are not the
same measurement --- none of the published metrics is this work's point-level car F$_1$,
and the published rows are all-class test-set scores while this work's is a car-only,
single-sequence one --- so the lower accuracy is a difference of paradigm, not a lower
score on a shared axis, and this work does not contest the gap. What the pipeline offers
instead is a different set of properties: it
segments without any labelled point clouds, it runs on a single consumer GPU without a
training stage for detection (its runtime is reported in Section~\ref{sec:results-runtime}), and every stage of
its decision is inspectable, so its
failure modes can be --- and are --- root-caused rather than left opaque (the recall
shortfall is traced to clustering fragmentation, not classifier error, in Section~\ref{sec:recall-analysis}).
The accuracy gap is therefore a deliberate consequence of choosing interpretability and
near-zero annotation over benchmark accuracy, not an unexplained deficiency.

A related question the table anticipates is why the comparison is drawn against
segmentation methods rather than 3D object detection, the task the ``detector'' label
most naturally invokes. Detection benchmarks score oriented 3D bounding boxes by average
precision, as in PointPillars~\cite{lang2019pointpillars} and
CenterPoint~\cite{yin2021centerpoint}. This pipeline, however, emits point-level clusters
--- it must, because those clusters are the input to the completion stage (Section~\ref{ch:completion-method}) ---
and never fits an oriented box; box AP would require a box-fitting stage the pipeline
deliberately does not have. Point-level IoU against the ground-truth car points is
therefore the metric that matches the actual output, and the choice of IoU threshold and
recall denominator that make that measurement precise --- together with its sensitivity
to those choices --- is examined directly in Section~\ref{sec:eval-protocol}.

This framing also sets the boundary of the thesis's contribution. Classical
ground-removal-plus-clustering detection --- including recent
clustering-plus-classifier detectors for LiDAR vehicles~\cite{lim2025longrange} ---
a lightweight multi-object tracking stage used only to filter transient detections
across frames (in the tradition of 3D multi-object trackers such as
AB3DMOT~\cite{weng2020ab3dmot}, and not a contribution of this work), and PCN-style
completion are established building blocks, reimplemented and tuned here; the
contribution is methodological: a validated
metric for completion quality on real LiDAR where reference-based Chamfer distance is
invalid (Section~\ref{sec:donor-validation}). Two diagnostic findings sit alongside it: the root-cause
analysis of the recall bottleneck together with the
chain of negative repair results (Section~\ref{sec:recall-analysis}), and the finding that synthetic
pretraining is redundant for the classifier at scale (Chapter~\ref{ch:evaluation}). The positioning table
is the backdrop against which the contribution is argued: this is an interpretable,
annotation-light system whose value is in what it explains and measures, not in topping a
benchmark it was never designed to win.

\section{Theoretical Background}
\label{sec:bg-theory}

This section fixes notation and states the standard definitions for the operators the
pipeline is built from, in the order a scan passes through them: ground-plane fitting, voxel
downsampling, density-based clustering, the point-set classifier, and the point-set
distances used for completion and its evaluation. Each is a textbook construction cited to
its source; the pipeline's own architectures and parameter choices are given in the method
chapters. A single LiDAR scan is written as a finite point set
$P = \{\mathbf{x}_i \in \mathbb{R}^3\}_{i=1}^{N}$.

\subsection{Ground-Plane Removal by RANSAC}
\label{sec:th-ransac}
A plane is a pair $(\mathbf{n}, d)$ with unit normal $\mathbf{n}\in\mathbb{R}^3$,
$\lVert\mathbf{n}\rVert = 1$, and offset $d$; a point lies on it when
$\mathbf{n}\cdot\mathbf{x} + d = 0$, and its signed distance to the plane is
\begin{equation}
  \delta(\mathbf{x};\mathbf{n},d) = \mathbf{n}\cdot\mathbf{x} + d .
\end{equation}
RANSAC~\cite{fischler1981random} fits the ground plane by repeatedly sampling three points,
forming the plane through them, and scoring it by its inlier set --- the points closer than
a tolerance $\tau_g$,
\begin{equation}
  \mathcal{I}(\mathbf{n},d) = \bigl\{\, \mathbf{x}\in P : \lvert \delta(\mathbf{x};\mathbf{n},d)\rvert \le \tau_g \,\bigr\}.
\end{equation}
The candidate with the largest inlier set is kept and its inliers removed as ground. To draw
at least one all-inlier sample with probability $p$, given inlier fraction $w$ and $s=3$
points per sample, the required number of iterations is
\begin{equation}
  T = \left\lceil \frac{\log(1-p)}{\log\!\bigl(1 - w^{s}\bigr)} \right\rceil .
\end{equation}

\subsection{Voxel-Grid Downsampling}
\label{sec:th-voxel}
To reduce the range-dependent density bias before clustering, the non-ground points are
downsampled on a regular grid of side $v$. Each point maps to an integer cell index
\begin{equation}
  \mathrm{vox}(\mathbf{x}) = \left\lfloor \frac{\mathbf{x} - \mathbf{x}_{\min}}{v} \right\rfloor \in \mathbb{Z}^3
\end{equation}
(the floor taken per coordinate), and every occupied cell $c$ is replaced by the centroid of
the points that fall in it,
\begin{equation}
  \bar{\mathbf{x}}_c = \frac{1}{\lvert P_c\rvert} \sum_{\mathbf{x}\in P_c} \mathbf{x} .
\end{equation}

\subsection{Density-Based Clustering: DBSCAN and HDBSCAN}
\label{sec:th-clustering}
The non-ground points are grouped into objects by density. DBSCAN~\cite{ester1996dbscan}
uses a fixed radius $\varepsilon$ and a minimum count $m$: with the
$\varepsilon$-neighbourhood
$\mathcal{N}_\varepsilon(\mathbf{p}) = \{\mathbf{q}\in P : \lVert\mathbf{p}-\mathbf{q}\rVert \le \varepsilon\}$,
a point is a \emph{core point} when
\begin{equation}
  \lvert \mathcal{N}_\varepsilon(\mathbf{p}) \rvert \ge m ,
\end{equation}
and clusters are the maximal sets connected through core points. A single $\varepsilon$
cannot serve both the dense near field and the sparse far field of a scan.
HDBSCAN~\cite{campello2013hdbscan,mcinnes2017hdbscan} removes the fixed radius. With the
\emph{core distance} $d_{\mathrm{core}}(\mathbf{p})$ set to the distance from $\mathbf{p}$ to
its $m$-th nearest neighbour, it defines the \emph{mutual reachability distance}
\begin{equation}
  d_{\mathrm{mreach}}(\mathbf{a},\mathbf{b}) = \max\!\bigl(\, d_{\mathrm{core}}(\mathbf{a}),\; d_{\mathrm{core}}(\mathbf{b}),\; \lVert\mathbf{a}-\mathbf{b}\rVert \,\bigr),
\end{equation}
builds a minimum spanning tree under it, and condenses the tree into a hierarchy indexed by
$\lambda = 1/d_{\mathrm{mreach}}$. A candidate cluster $C$ born at $\lambda_{\min}(C)$, losing
each point $\mathbf{p}$ at level $\lambda_{\mathbf{p}}$, has stability
\begin{equation}
  S(C) = \sum_{\mathbf{p}\in C} \bigl(\lambda_{\mathbf{p}} - \lambda_{\min}(C)\bigr),
\end{equation}
and the flat clustering is the set of disjoint clusters of maximum total stability.
Stability-based extraction is what lets one setting cluster both near and far cars; it is
also, as Section~\ref{sec:recall-analysis} shows, what splits a single car when an internal
density gap makes two sub-clusters individually more stable than the whole.

\subsection{Permutation-Invariant Point Classification}
\label{sec:th-pointnet}
The car/not-car classifier is built on PointNet~\cite{Qi_2017_CVPR}, whose defining property
is invariance to the ordering of its input points. A function of a point set is made
permutation-invariant by applying a shared map $h$ to each point, pooling with a symmetric
operator, and passing the pooled feature through a final map $g$:
\begin{equation}
  f(\{\mathbf{x}_1,\dots,\mathbf{x}_n\}) = g\!\left( \max_{i=1,\dots,n} h(\mathbf{x}_i) \right),
\end{equation}
the maximum taken coordinate-wise over the per-point feature vectors. Because the $\max$ pool
is symmetric, $f$ is independent of the order of the points; $h$ and $g$ are multilayer
perceptrons learned end-to-end. Chapter~\ref{ch:pipeline} applies this design to a single
cluster.

\subsection{Point-Set Distances: Chamfer, IoU, and Donor Coverage}
\label{sec:th-distances}
Completion and its evaluation are expressed through distances between point sets. The
\emph{Chamfer distance} between sets $A$ and $B$ averages, in each direction, the squared
distance from every point to the nearest point of the other set,
\begin{equation}
  \mathrm{CD}(A,B) = \frac{1}{\lvert A\rvert}\sum_{\mathbf{a}\in A} \min_{\mathbf{b}\in B} \lVert\mathbf{a}-\mathbf{b}\rVert^2 \;+\; \frac{1}{\lvert B\rvert}\sum_{\mathbf{b}\in B} \min_{\mathbf{a}\in A} \lVert\mathbf{b}-\mathbf{a}\rVert^2 ,
\end{equation}
and is the standard training objective for completion on synthetic
shapes~\cite{fan2017point,yuan2019pcnpointcompletionnetwork};
Section~\ref{sec:donor-validation} shows it is invalid as a real-data \emph{score}. Detection is
scored by a point-level intersection-over-union between a predicted cluster $Q$ and a
ground-truth car's point set $G$,
\begin{equation}
  \mathrm{IoU}(Q,G) = \frac{\lvert Q \cap G\rvert}{\lvert Q \cup G\rvert},
\end{equation}
a prediction matching a ground-truth object when $\mathrm{IoU}(Q,G)\ge 0.3$
(Chapter~\ref{ch:evaluation}). Completion is scored by the \emph{donor coverage} developed
in Section~\ref{sec:donor-validation}. Given the input cluster $I$, a \emph{donor cloud} $D$ of
the same car's labelled points accumulated from other frames, and a visibility radius
$\tau$, the \emph{novel set} is the donor surface the input never observed,
\begin{equation}
  \mathcal{V} = \bigl\{\, \mathbf{d}\in D : \min_{\mathbf{p}\in I} \lVert\mathbf{d}-\mathbf{p}\rVert \ge \tau \,\bigr\},
\end{equation}
and the coverage of a candidate cloud $C$ at radius $r$ is the fraction of the novel set it
reaches,
\begin{equation}
  \mathrm{cov}_r(C) = \frac{1}{\lvert\mathcal{V}\rvert} \Bigl\lvert \bigl\{\, \mathbf{d}\in\mathcal{V} : \min_{\mathbf{c}\in C} \lVert\mathbf{d}-\mathbf{c}\rVert < r \,\bigr\} \Bigr\rvert .
\end{equation}
By construction the raw partial covers none of $\mathcal{V}$; the full construction, its
mirrored baseline, and its per-band hallucination guard are developed in
Section~\ref{sec:donor-validation}.

Three points emerge from this review and set up the rest of the thesis. First, the accuracy
of mainstream LiDAR segmentation is inseparable from a dense per-point annotation budget,
which motivates the interpretable, near-annotation-free clustering-plus-classifier route
taken here --- and whose recall cost the later chapters measure and root-cause. Second, that
route's stage-level transparency is what makes such root-causing possible at all, where an
end-to-end network offers no stage to inspect. Third, and most consequentially for the
contribution, completion on real LiDAR has no valid reference-based score: every established
metric is anchored to the observed input, a synthetic prior, or a one-sided accumulation, so
none measures recovered unseen surface. These three points are the basis on which the thesis
is built --- the last is the gap its two research questions address
(Chapter~\ref{ch:introduction}), by defining and validating a donor-frame coverage metric and
using it to test whether completion adds genuine unseen surface and improves downstream boxes.
The chapter has also fixed the notation and standard operators the later chapters build on.
Chapter~\ref{ch:pipeline} turns those operators into the detection pipeline and the completion
method, and develops the donor-frame coverage metric sketched here in full.
